\documentclass[10pt,a4paper]{amsart}
\usepackage{amsmath,amssymb,amsthm,graphicx}
% Retained minimal preamble for the Stacks Project chapter sites-cohomology.tex (Cohomology on Sites).
% Verbatim excerpt of the chapter preamble: Xy-pic setup (native preamble lines 14-18) and the native
% theorem declarations with the mathematical macros (native preamble lines 154-217). Copied mechanically;
% nothing added, removed or redefined. The wrapper class and the amsmath/amssymb/amsthm/graphicx base
% packages are supplied by the assembler, so the class/externaldocument part of the native preamble is omitted.
% For commutative diagrams we use Xy-pic
\usepackage[all]{xy}

% We use 2cell for 2-commutative diagrams.
\xyoption{2cell}
% Theorem environments.
%
\theoremstyle{plain}
\newtheorem{theorem}[subsection]{Theorem}
\newtheorem{proposition}[subsection]{Proposition}
\newtheorem{lemma}[subsection]{Lemma}

\theoremstyle{definition}
\newtheorem{definition}[subsection]{Definition}
\newtheorem{example}[subsection]{Example}
\newtheorem{exercise}[subsection]{Exercise}
\newtheorem{situation}[subsection]{Situation}

\theoremstyle{remark}
\newtheorem{remark}[subsection]{Remark}
\newtheorem{remarks}[subsection]{Remarks}

\numberwithin{equation}{subsection}

% Macros
%
\def\lim{\mathop{\mathrm{lim}}\nolimits}
\def\colim{\mathop{\mathrm{colim}}\nolimits}
\def\Spec{\mathop{\mathrm{Spec}}}
\def\Hom{\mathop{\mathrm{Hom}}\nolimits}
\def\Ext{\mathop{\mathrm{Ext}}\nolimits}
\def\SheafHom{\mathop{\mathcal{H}\!\mathit{om}}\nolimits}
\def\SheafExt{\mathop{\mathcal{E}\!\mathit{xt}}\nolimits}
\def\Sch{\mathit{Sch}}
\def\Mor{\mathop{\mathrm{Mor}}\nolimits}
\def\Ob{\mathop{\mathrm{Ob}}\nolimits}
\def\Sh{\mathop{\mathit{Sh}}\nolimits}
\def\NL{\mathop{N\!L}\nolimits}
\def\CH{\mathop{\mathrm{CH}}\nolimits}
\def\proetale{{pro\text{-}\acute{e}tale}}
\def\etale{{\acute{e}tale}}
\def\QCoh{\mathit{QCoh}}
\def\Ker{\mathop{\mathrm{Ker}}}
\def\Im{\mathop{\mathrm{Im}}}
\def\Coker{\mathop{\mathrm{Coker}}}
\def\Coim{\mathop{\mathrm{Coim}}}

% Boxtimes
%
\DeclareMathSymbol{\boxtimes}{\mathbin}{AMSa}{"02}

%
% Macros for moduli stacks/spaces
%
\def\QCohstack{\mathcal{QC}\!\mathit{oh}}
\def\Cohstack{\mathcal{C}\!\mathit{oh}}
\def\Spacesstack{\mathcal{S}\!\mathit{paces}}
\def\Quotfunctor{\mathrm{Quot}}
\def\Hilbfunctor{\mathrm{Hilb}}
\def\Curvesstack{\mathcal{C}\!\mathit{urves}}
\def\Polarizedstack{\mathcal{P}\!\mathit{olarized}}
\def\Complexesstack{\mathcal{C}\!\mathit{omplexes}}
% \Pic is the operator that assigns to X its picard group, usage \Pic(X)
% \Picardstack_{X/B} denotes the Picard stack of X over B
% \Picardfunctor_{X/B} denotes the Picard functor of X over B
\def\Pic{\mathop{\mathrm{Pic}}\nolimits}
\def\Picardstack{\mathcal{P}\!\mathit{ic}}
\def\Picardfunctor{\mathrm{Pic}}
\def\Deformationcategory{\mathcal{D}\!\mathit{ef}}
\usepackage{cleveref}
\crefname{originalsection}{original section}{original sections}
\Crefname{originalsection}{Original section}{Original sections}
\makeatletter
\newlabel{sites-section-sheaves-injective}{{00WL}{source}{}{sourcepointer-1}{}}
\newlabel{sites-section-sheaves-injective@cref}{{[originalsection][0][]00WL}{[1][0][]source}{}{}{}}
\makeatother
\title{Lifting a map from a summand of a finite free module along a surjection, locally}
\author{The Stacks Project Authors}
\date{stacks-project:sites-cohomology.tex:local-snapshot-20261011; GNU FDL 1.2 or later}
\begin{document}
\maketitle
\noindent\textit{Editorial note.} All delivered mathematical text is the author's own, copied from the pinned Stacks Project chapter file sites-cohomology.tex; no statement or proof step was rewritten, reordered or removed.\par\medskip
\noindent\textit{Editorial note.} The proof is the author's complete argument: the reduction to the free case, the translation into lifted sections and the local liftability are all inside the counted span.\par\medskip
\noindent\textit{Editorial note.} Cross-chapter restatement: the same label occurs in chapter cohomology for a ringed space. That statement is the special case of this ringed-site statement, so per the lane-s4 rule only this more general occurrence is counted here.\par\medskip
\noindent\textit{Editorial note.} The single labelled reference occurring inside the counted spans points to a chapter of the Project not staged in this card's source tree; it is carried in assembly.external\_references as a non-load-bearing navigation pointer recorded at the invoking line.\par\medskip
\noindent\textit{Editorial note.} The preamble is the chapter's own retained minimal preamble (sources/stacks-sites-cohomology/retained-preamble.tex), a verbatim excerpt of the native preamble: the Xy-pic setup and the native theorem declarations with the mathematical macros.\par\medskip
\noindent\textit{Allowed standard prerequisites (editorial).}\begin{itemize}
\item Direct summand of a finite free module on a ringed site. Same chapter, the notion of a direct summand of a finite free O\_U-module, which is the concrete form of a finite projective module. Supplies the hypothesis on E, and the reduction to the case E = O\_U\textasciicircum{}\{direct sum n\}.
\item Characterisation of surjective maps of sheaves. Sites, the characterisation of an epimorphism of sheaves by local liftability of sections. Supplies the step that the lift of the images of the basis elements is possible on the members of a covering.
\item Sections of a finite free module and their description by basis elements. Same chapter, the identification of Hom(O\_U\textasciicircum{}\{direct sum n\}, F) with tuples of sections of F over U. Supplies the equivalence between a dotted arrow and a lift of the images of the basis elements.
\end{itemize}\medskip
\ifdefined\hypertarget\hypertarget{sourcepointer-1}{}\fi
\noindent\textit{Source reference.} Original section 00WL: sources/stacks-sites-cohomology/pinned-original-source.tex, source lines 11792--11792. Cross-chapter navigation pointer invoked at proof lines 11791-11792, where the author uses the characterisation of surjective maps of sheaves. The target chapter is not staged in this card's source tree, so the pointer records the invoking line; it only resolves the reference for the reader, and no bibliography citation is involved.\par\medskip
\begin{lemma}
\label{lemma-local-lift-map}
Let $(\mathcal{C}, \mathcal{O})$ be a ringed site. Let $U$ be an object of
$\mathcal{C}$. Given a solid diagram of $\mathcal{O}_U$-modules
$$
\xymatrix{
\mathcal{E} \ar@{..>}[dr] \ar[r] & \mathcal{F} \\
& \mathcal{G} \ar[u]_p
}
$$
with $\mathcal{E}$ a direct summand of a finite free
$\mathcal{O}_U$-module and $p$ surjective, then there exists a
covering $\{U_i \to U\}$ such that a dotted arrow
making the diagram commute exists over each $U_i$.
\end{lemma}

\begin{proof}
We may assume $\mathcal{E} = \mathcal{O}_U^{\oplus n}$ for some $n$.
In this case finding the dotted arrow is equivalent to lifting the
images of the basis elements in $\Gamma(U, \mathcal{F})$. This is
locally possible by the characterization of surjective maps of
sheaves (Sites, Section \ref{sites-section-sheaves-injective}).
\end{proof}

\end{document}
